% =====================================================================
%  ESEMPIO di verbale interno: mostra tutto ciò che il modello sa fare.
%  Il contenuto è di prova: sostituitelo con quello della vostra
%  riunione reale prima di consegnarlo.
% =====================================================================
\documentclass[verbale=interno]{gpt}

\titolo{Analisi dei capitolati}
\data{2026-10-15}
\orario{17:30}{19:00}
\luogo{Discord}
\redattore{guglielmo}
\verificatore{biasini}
\componenti{biasini,cescobolla,corro,guglielmo,mazzurana,spagnoli}
\assenti{spagnoli}

\modifica{1.0.0}{2026-10-16}{guglielmo}{biasini}{Verifica e approvazione}
\modifica{0.1.0}{2026-10-15}{guglielmo}{--}{Prima stesura}

\begin{document}

\section{Ordine del giorno}
\begin{enumerate}
  \item Confronto tra i capitolati presentati.
  \item Scelta dei capitolati di interesse.
  \item Domande da inviare ai proponenti.
  \item Strumenti di lavoro del gruppo.
\end{enumerate}

\section{Svolgimento}

\subsection{Confronto tra i capitolati}
Ogni membro ha letto i dieci capitolati. Il gruppo ha valutato la
fattibilità di ciascuno considerando il numero di tecnologie nuove,
la presenza di hardware, la chiarezza del perimetro e il supporto
offerto dal proponente. Sono stati scartati i capitolati con hardware
fisico (C3) e quelli con uno stack vincolato molto ampio (C5, C6, C8).

\subsection{Scelta dei capitolati di interesse}
Restano in valutazione C0 (AIRA), C1 (AURA) e C2 (NEXUM Concierge).
Per C1 il proponente fornisce un kit di impianto fisico e l'intero
perimetro del progetto è il server \glo{MCP} riportato in
Figura~\ref{fig:aura}.

\figura[0.9]{esempi/2026-10-15/architettura-aura.png}{Architettura di
riferimento di AURA (fonte: presentazione di Vimar S.p.A.)\label{fig:aura}}

\begin{nota}[Criticità emerse]
Per C0 resta da chiarire se le richieste vadano validate da un operatore
o inviate in automatico; per C2 non è ancora disponibile il capitolato
completo.
\end{nota}

\subsection{Strumenti di lavoro}
Il gruppo userà GitHub per codice e documentazione, Discord per le
riunioni e questo modello \LaTeX{} per tutti i verbali.

\section{Decisioni}
\begin{decisioni}
  \decisione{I capitolati di interesse sono C0, C1 e C2, in ordine di
    preferenza da stabilire dopo le risposte dei proponenti.}
  \decisione{I verbali si redigono con il modello \LaTeX{} comune e si
    archiviano nel repository di documentazione.}
  \decisione{Ogni verbale ha un redattore e un verificatore diversi.}
\end{decisioni}

\section{Azioni da intraprendere}
\begin{azioni}
  \azione[\#1]{Inviare le domande a Vimar (Sciacco, in copia
    Bevilacqua).}{guglielmo}{17/10/2026}
  \azione[\#2]{Chiedere a Sanmarco Informatica chiarimenti sulla
    validazione delle richieste.}{corro}{17/10/2026}
  \azione[\#3]{Recuperare il capitolato completo di C2.}{mazzurana}{18/10/2026}
  \azione[\#4]{Creare l'organizzazione GitHub e il repository dei
    documenti.}{cescobolla,biasini}{19/10/2026}
\end{azioni}

\section{Prossima riunione}
Lunedì 19 ottobre 2026 alle 17:30 su Discord, per discutere le risposte
dei proponenti e stilare la classifica dei capitolati.

\end{document}
